The Indiana Network for Patient Care (INPC) harmonizes and combines electronic medical records (EMR) from multiple health systems, hospitals, and clinics in Indiana \citep{regenstrief2023}. The EMR data contain patient demographics and healthcare encounter records that occur at contributing health care providers, including diagnosis codes, laboratory tests, and laboratory test results. Patients can be followed longitudinally using an encrypted person identifier. \footnote{All data cleaning and analyses were conducted using R version 4.4.1, and ChatGPT and Claude were used for code debugging. }

Our study sample consists of infants born at INPC facilities between 2019 and 2023 who can be linked to a birth mother. We exclude infant-mother pairs with a residential zip code outside of Indiana, maternal delivery age less than 12 or older than 60, missing maternal gender or maternal gender coded as male, or a discrepancy between infant date of birth and recorded date of delivery. We limit the sample to infants with race coded as White or Black due to low cell sizes among the other racial/ethnic groups in our data.

Our final analytic sample contains one record for each of the $N = 140,562$ newborn infants who met inclusion criteria. The main outcome is a binary variable indicating whether a laboratory test for in-utero drug exposure was conducted on the infant. We include laboratory tests for opioids, benzodiazepines, cocaine, methamphetamine, amphetamines, and other drugs of abuse. Specifically, an infant is classified as tested if their medical record contains a lab test with a LOINC code corresponding to a family of meconium or umbilical cord tests for drugs of abuse -- the full list of codes is in Appendix Table \ref{loinc_codes}. A secondary outcome, measured only for infants who were tested, is a binary indicator that the test was positive for in-utero drug exposure. 

In addition to lab test data, we also constructed a rich set of baseline covariates for each infant-mother pair. For the mother, we measure delivery age; the number of prenatal care visits; racial group; maternal insurance status (Medicaid, private, or other); and health history including prior diagnosis of chronic pain disorder, gestational diabetes, hypertensive disorder, mental health disorders, or substance use disorder. For the infant, we measure diagnosis of premature birth, low birth weight, and the number of healthcare encounters within one month of birth, which we view as a summary measure of infant health status. 

We also attached two geographic measures of the severity of the opioid epidemic in the infant's home county. First, we used Centers for Disease Control and Prevention (CDC) National Vital Statistics Mortality to compute the crude drug overdose mortality rate from 2014-2018 per 10,000 for each county in Indiana. Second, we used Indiana inpatient data from the Agency for Healthcare Research and Quality (AHRQ) Healthcare Cost and Utilization Project (HCUP) from 2018 to compute the rate of neonatal abstinence syndrome diagnosis per 100 births for each county in Indiana. We used the infant's home address to attach this county level information to each infant in our sample.
